\section*{ЗАКЛЮЧЕНИЕ}
\addcontentsline{toc}{section}{ЗАКЛЮЧЕНИЕ}

Выпускная квалификационная работа была посвящена разработке и внедрению веб-платформы для автоматизации бизнес-процессов управления персоналом компании. Актуальность проекта обусловлена необходимостью повышения эффективности внутренних коммуникаций, устранения фрагментированности цифровых сервисов и централизации управления задачами и данными в рамках одной системы.

В ходе выполнения проекта были достигнуты следующие результаты:

\begin{enumerate}
\item Проведен анализ предметной области, в результате которого выявлены ключевые проблемы коммуникации и управления в организации. Выбрана архитектура и технологии, соответствующие задачам проекта.
\item Разработана концептуальная модель веб-платформы, определены функциональные модули и их взаимодействие. Проект охватывает такие компоненты, как почта, календарь, чаты, управление проектами, видеоконференции, хранилище файлов и контактная книга.
\item Осуществлено проектирование пользовательского интерфейса и архитектуры платформы с учетом требований к масштабируемости, безопасности и адаптивности.
\item Реализована серверная и клиентская части системы с применением современных технологий (React, TypeScript, Node.js, JMAP, WebSocket, PostgreSQL).
\item Проведено модульное и системное тестирование платформы, внедрены механизмы безопасности и мониторинга, осуществлено развертывание системы в продуктивной среде.
\end{enumerate}

Разработанная веб-платформа была успешно внедрена в производственной среде компании АО «Белуга Проджектс Лоджистик», где доказала свою эффективность в решении поставленных задач. Все цели, сформулированные в начале работы, были достигнуты, а технические и функциональные требования — реализованы в полном объеме.

Проект демонстрирует высокий потенциал для дальнейшего масштабирования и развития, а также может быть адаптирован под нужды других предприятий, нуждающихся в комплексной системе управления персоналом и коммуникацией.